% =====================================================================
% 4. Implementation
% =====================================================================
\section{Implementation}
\label{sec:impl}

In this section, we discuss the pipeline for the generation of
self-captured data, personalization process consisting of seven
steps, the network architecture, the training process and two
technological innovations of the current research project --
subject-aware sampling and two-stage personalization process.

\subsection{System Pipeline}
\label{sec:impl:pipeline}

As seen in Figure~\ref{fig:pipeline}, there are seven stages of
the pipeline. Each stage is a Python script which is located
either in the project root folder or in the subfolder
\texttt{training}.

\begin{figure}[H]
  \centering
  \includegraphics[width=\linewidth]{flowchart.png}
  \vspace{-18pt}
  \caption{\footnotesize End-to-end pipeline. Blue boxes are data preparation
  stages; orange boxes are the two stages contributed by this
  project; green boxes are the saved model files that feed the
  live demo.}
  \label{fig:pipeline}
\end{figure}

Steps one to four collect self-captured data. The script named
\texttt{capture.py} captures the video stream from the laptop
webcam and writes the captured frames to the file in the
corresponding directories depending on the subject and session
number. The script \texttt{prelabel.py} produces weak labels
with the help of a small pre-trained model so that the reviewer
had something to start with. The script \texttt{review\_app.py}
is a small web app based on Streamlit framework for a team
member to go through all the frames and correct or confirm the
label. The script \texttt{gather.py} aggregates the individual
subject CSV files to \texttt{prelabels.csv} file in the project
root directory taking care of duplicate elimination and path
normalization so that the frames captured on Windows system have
the same path format as the frames captured on a macOS system.

The remaining three steps involve preparation of the training
set and training of the network. The script
\texttt{face\_crop.py} applies MTCNN on all frames from
self-captured dataset and saves the tight crop of
$224 \times 224$ size in a separate directory tree
(\texttt{raw\_face/}). The script \texttt{merge\_ferplus.py}
merges the self-captured subset with FER+ and RAF-DB datasets
into three CSV files (\texttt{dataset/train.csv},
\texttt{dataset/val.csv} and \texttt{dataset/test.csv}). The
script \texttt{train.py} loads these files, applies
subject-aware sampler to the dataset and trains the ResNet-18
network with saving the best checkpoint by validation accuracy.
Finally, the script \texttt{finetune.py} loads generic
checkpoint, makes personalization step on a target subject and
saves the personalized checkpoint.

\subsection{Model Architecture and Backbone Selection}
\label{sec:impl:arch}

ResNet-18 was chosen as the backbone after reviewing three
different architectures that are widely used for FER. The VGG-16
architecture \autocite{simonyan2015vgg} consists of a chain of
$3 \times 3$ convolutions and is able to achieve high accuracy
in FER tasks, but its large number of $138.4$ million parameters
results in long training and deployment times on the laptop
where the application is hosted. EfficientNet-B0
\autocite{tan2019efficientnet}, on the contrary, is the most
parameter-efficient architecture with $4.0$ million parameters,
but in the comparison presented in
Section~\ref{sec:results:archabl} it showed somewhat worse
results on all slices compared to ResNet-18. ResNet-18
architecture lies between those two with $11.2$ million
parameters, as it utilizes residual connections
\autocite{he2016resnet} to avoid the decrease in accuracy that
is inherent to deeper plain networks, it can be fine-tuned with
ImageNet weights due to early layers' ability to recognize
common features like edges and textures that are relevant for
any image classification task, and it is fit in one laptop GPU
with the batch size of $128$. Thus, we choose ResNet-18
architecture for our task and report the comparison with the
VGG-16 in Section~\ref{sec:results:archabl}.

The proposed model is based on ResNet-18 backbone trained on
ImageNet dataset with pretrained weights, but instead of a head
consisting of $1000$ class classifiers there is a linear layer
projecting $512$-dimensional pooled feature into our emotions
space. Figure~\ref{fig:arch} shows the structure of the network
and marks the trainable blocks on each stage.

\begin{figure}[H]
  \centering
  \resizebox{\linewidth}{!}{%
  \begin{tikzpicture}[node distance=3.5mm]
    \tikzstyle{blk}=[gap2block, minimum width=22mm, minimum height=8mm,
                     font=\footnotesize\bfseries, fill=blue!7]
    \tikzstyle{frz}=[blk, fill=gray!8, draw=gray!60]
    \tikzstyle{tun}=[blk, fill=orange!15, draw=orange!70!black]
    \tikzstyle{io}=[blk, fill=green!10, draw=green!60!black]

    \node[io]                    (in)  {$224\times224$ RGB};
    \node[frz, right=of in]      (s1)  {conv1+bn1};
    \node[frz, right=of s1]      (s2)  {layer1};
    \node[frz, right=of s2]      (s3)  {layer2};
    \node[frz, right=of s3]      (s4)  {layer3};
    \node[tun, right=of s4]      (s5)  {layer4};
    \node[tun, right=of s5]      (s6)  {avgpool+fc};
    \node[io,  right=of s6]      (out) {$7$ logits};

    \draw[gap2arr] (in) -- (s1);
    \draw[gap2arr] (s1) -- (s2);
    \draw[gap2arr] (s2) -- (s3);
    \draw[gap2arr] (s3) -- (s4);
    \draw[gap2arr] (s4) -- (s5);
    \draw[gap2arr] (s5) -- (s6);
    \draw[gap2arr] (s6) -- (out);
  \end{tikzpicture}}
  \caption{\footnotesize ResNet-18 backbone with a $7$-class linear head. At
  generic training all blocks are trainable. At the
  personalization stage the gray blocks
  (\texttt{conv1}, \texttt{bn1}, \texttt{layer1},
  \texttt{layer2}, \texttt{layer3}) are frozen and only the
  orange blocks (\texttt{layer4} and the head) receive gradient
  updates, which corresponds to $8{,}397{,}319$ trainable
  parameters out of $11{,}180{,}103$ in total ($75.1\%$).}
  \label{fig:arch}
\end{figure}

Each input is scaled to size $224 \times 224$ and normalized
using the mean and standard deviation used in ImageNet
normalization. In training, we also add random crop, random
horizontal flip, color jitter, and random erasing as the common
data augmentations. For the images in the self-capture videos,
we use the directory containing faces cropped with MTCNN if it
exists; in FER+ and RAF-DB, the faces have been pre-cropped and
aligned in the corresponding dataset creation scripts. At
evaluation time, we also perform horizontal flip test-time
augmentation: we average the softmax probabilities of the
original image and the horizontal mirroring of it and then take
the argmax.

\subsection{Training Procedure}
\label{sec:impl:training}

The generic network is trained using cross-entropy loss over the
seven classes, with weights computed to balance class
distribution and label smoothing factor $0.1$
\autocite{szegedy2016rethinking}. Weights for each class are
calculated according to the following formula:
$w_c = N / (C \cdot N_c)$, where $N$ is the number of training
frames, $C = 7$ is the number of classes and $N_c$ is the count
of class $c$. After user2 contribution, classes of disgust and
fear still get the biggest weights (around $5.4$--$5.9$), since
they are still the most rare classes after discarding the
contempt class.

The optimizer we use is AdamW \autocite{loshchilov2019adamw}
with learning rate $1 \times 10^{-4}$ and weight decay
$1 \times 10^{-4}$. The learning rate is reduced from the
starting value using cosine schedule, reaching $0$ over the
planned $30$ epochs of training. We use a batch size of $128$ on
NVIDIA T4 GPU. The early stopping procedure is used, with
patience set to $8$ epochs; however, in the final run, early
stopping was never triggered, so training ran to full $30$
epochs. Best checkpoint with the highest validation accuracy is
saved to \texttt{checkpoints/best.pt}. All hyperparameters are
summarized in Table~\ref{tab:hparam} below.

\begin{table}[H]
  \centering
  \footnotesize
  \begin{tabular}{lll}
    \toprule
    Group & Hyperparameter & Value \\
    \midrule
    Optimizer & AdamW \autocite{loshchilov2019adamw} & $\beta_1=0.9$, $\beta_2=0.999$ \\
              & Learning rate         & $1\times 10^{-4}$ \\
              & Weight decay          & $1\times 10^{-4}$ \\
              & LR schedule           & Cosine annealing, $T_{\max}=30$ \\
    Loss      & Cross-entropy weights & Inverse class frequency \\
              & Label smoothing       & $0.1$ \\
    Sampling  & Self oversample       & $20\times$ (Sec.~\ref{sec:impl:sampling}) \\
    Augment.  & Resize $\rightarrow$ crop & $256 \rightarrow 224$ \\
              & Random horizontal flip & $p=0.5$ \\
              & Colour jitter          & $0.2$ each channel \\
              & Random erasing         & $p=0.1$, scale $(0.02,0.1)$ \\
              & mixup \autocite{zhang2018mixup} & $\alpha=0.2$ \\
              & Test-time augmentation & H-flip averaging \\
    Schedule  & Batch size            & $128$ \\
              & Epochs                & $30$ \\
              & Patience              & $8$ (early stop on val acc) \\
              & Image size            & $224\times 224$, RGB \\
    Hardware  & Accelerator           & NVIDIA T4 (GCP) \\
    \bottomrule
  \end{tabular}
  \caption{\footnotesize Hyperparameters used to train the generic model.
  Gradient-based optimization and weighted random sampling are
  topics covered in the course; mixup, label smoothing and
  test-time augmentation are common regularizers documented in
  the cited papers.}
  \label{tab:hparam}
\end{table}

Figure~\ref{fig:curves} shows the per-epoch values of training
and validation loss and per-source validation accuracy. The
training loss steadily decreases from about $1.46$ in epoch $1$
to $0.66$ in epoch $30$ without any late-stage divergence.
Validation loss closely follows the training loss but does not
keep falling beyond epoch $20$, which corresponds to plateauing
of the overall validation accuracy around this epoch as well.
Per-source graphs on the right panel show how the cross-domain
gap becomes evident during training: the validation accuracies
of FER+ and RAF-DB increase to above $80\%$ already in the first
$10$ epochs, while validation accuracy of user1 saturates at
around $58\%$ even with the subject-aware oversampler on. This
is the empirical signature that motivates the personalization
stage of Section~\ref{sec:impl:personalize}.

\begin{figure}[H]
  \centering
  \includegraphics[width=\linewidth]{figures/training_curves.png}
  \caption{\footnotesize Per-epoch training curves for the generic
  ResNet-18 model. Left: training and validation
  cross-entropy loss decrease together over $30$ epochs, with
  the training loss falling from $1.46$ to $0.66$ and the
  validation loss plateauing around epoch $20$ at $0.93$.
  Right: validation accuracy decomposed by source. The
  public-source slices (FER+, RAF-DB) climb fastest; the
  self-captured slice plateaus much lower, foreshadowing the
  cross-subject failure analyzed in
  Section~\ref{sec:results:failure}.}
  \label{fig:curves}
\end{figure}

\subsection{Subject-Aware Sampling}
\label{sec:impl:sampling}

From Table~\ref{tab:datasummary}, it can be seen that
self-captured slice makes up merely $1.4\%$ of the whole
training set. Uniform sampling of batch size $128$ samples will
provide less than two self-captured samples on average. This is
insufficient for the backbone to learn the kind of frames it
will see at inference time.

To solve this problem, we use the
\texttt{WeightedRandomSampler} as an argument of the PyTorch
data loader where each self-captured sample gets weight $20$,
while every public-source sample retains weight $1$. The sampler
draws indices with replacement, so the expected number of
self-captured frames per epoch is about
$555 \times 20 = 11{,}100$, which is on the same order as the
public slice. The class balance of each public source is kept
the same because we set all public-source weights to $1$. The
sampler only rebalances along the source axis, not along the
class axis. Algorithm~\ref{alg:sampler} gives the pseudocode.
The actual implementation is a few lines of Python and uses the
standard PyTorch \texttt{DataLoader} API through one keyword
argument.

\begin{algorithm}[H]
\small
\caption{\footnotesize Subject-aware weighted random sampler}
\label{alg:sampler}
\begin{algorithmic}[1]
\Require Training set $\mathcal{D}_{\text{tr}}$ with per-row source
         label $s_i \in \{\text{self},\text{ferplus},\text{rafdb}\}$;
         oversample factor $\rho$
\Function{BuildSampler}{$\mathcal{D}_{\text{tr}}, \rho$}
  \State $\mathbf{w} \gets$ empty vector of length $|\mathcal{D}_{\text{tr}}|$
  \For{$i = 1$ to $|\mathcal{D}_{\text{tr}}|$}
    \State $w_i \gets \rho$ \textbf{if} $s_i = \text{self}$ \textbf{else} $1$
  \EndFor
  \State \Return \texttt{WeightedRandomSampler}$(\mathbf{w},
         \text{num\_samples}=|\mathcal{D}_{\text{tr}}|,
         \text{replacement}=\text{True})$
\EndFunction
\end{algorithmic}
\end{algorithm}

This technique is similar to the class oversampling strategy
applied in long-tailed image classification, except for the fact
that we address the data source imbalance rather than the class
imbalance. Increasing the value of $\rho$ to $20$ increases
validation accuracy of self-captured videos by five percentage
points and shifts validation accuracies of public sources by
less than one percentage point (see
Section~\ref{sec:results:overall}).

\subsection{Optional Personalization Layer}
\label{sec:impl:personalize}

Despite subject-aware sampling, the generic model achieves an
accuracy of merely $21.8\%$ on the unseen subject
(Section~\ref{sec:results:overall}). In order to bridge the gap,
we include an additional step during training to adapt the
network to a target user by using only a few images from the
target's own dataset.

The personalization step takes the \texttt{best.pt} model and
the target subject $u^\star$, and then reads all the frames of
$u^\star$ labeled by \texttt{prelabels.csv}, divides the frames
inside the subject (stratified $80/20$ split) to fine-tuning and
held-out test sets. Then, we keep frozen the layers
\texttt{conv1}, \texttt{bn1}, \texttt{layer1}, \texttt{layer2}
and \texttt{layer3} and train only \texttt{layer4} and the
\texttt{fc} ($8.40$M parameters out of $11.18$M, $75\%$). The
trainable parameters are optimized for $5$ epochs using AdamW,
with $lr = 1 \times 10^{-4}$, weight decay $1 \times 10^{-4}$,
batch size $32$ and plain cross-entropy loss (we do not use
class-weighting in this stage). $lr = 1 \times 10^{-4}$ is the
result of a small learning rate sweep that can be found in
Section~\ref{sec:results:personalize}
(Table~\ref{tab:ftlr}). Then, we evaluate both generic and
personalized checkpoints on the same $20\%$ held-out test. The
motivation behind freezing the first three residual stages comes
from a well-known result stating that early convolutional blocks
learn general features (like edges and textures), which transfer
well between users, whereas upper blocks and the network's head
learn subject-specific features
\autocite{finn2017maml, ganin2015dann}. The reason for not
freezing \texttt{layer4} is that almost $25\%$ of the network's
parameters belong to this layer. We also noticed that when
freezing \texttt{layer4}, the network under-fits the target
subject and the personalization effect becomes worse.

\subsection{Application Deployment as EmotionLens}
\label{sec:impl:deploy}

The personalized checkpoint serves as the inference engine
behind \emph{EmotionLens}, an interactive emotion recognition
web application. The backend consists of a FastAPI server, which
includes the WebSocket endpoint. The webcam frames are streamed
from the client to the server, where the MTCNN
\autocite{zhang2016mtcnn} detects faces; the detected face is
resized to $224 \times 224$ pixels and normalized using ImageNet
statistics; and the personalized ResNet-18 provides the
seven-class softmax output. Exponential moving average smoothing
is applied to the softmax vector to stabilize the predicted
class and prevent it from flickering frame by frame. The entire
process takes less than $50$~ms per frame on a modern laptop
GPU, making the overlays update at near-camera frame rate.

The application supports five separate analysis lenses built on
top of the same inference engine: the crowd emotion histogram,
the emotion alert, which triggers when negative emotions
persist, the emotion timeline, the public speaking nervousness
index, and the mimic game. Each of these lenses is represented
by a separate Python script, which analyzes the
seven-dimensional softmax, without requiring changes to the
original model. The frontend consists of HTML/CSS/JS with lens
switching performed by Swiper.js, and the web application is
available through a Cloudflare Tunnel, since browser
\texttt{getUserMedia} requires HTTPS.
